% !TEX root = ../main.tex
% !TEX program = lualatex
\documentclass[../main.tex]{subfiles}
\IfSubfilesClassLoaded{\externaldocument{\subfix{../../build/main}}}{}
% =============================================================================
% File: 23_Waterfront_Organization_Work_Control_Safety.tex
% =============================================================================
\begin{document}
\IfSubfilesClassLoaded{\chapter{EDO Study Guide}}{}
\section{Waterfront Organizations, Work Controls, and Safety}

\subsection{Naval Propulsion Program Oversight}
\begin{description}
  \item[\textbf{\ac{navreact}.}] Established by executive order and statute, Naval Reactors owns cradle-to-grave authority for nuclear propulsion design, maintenance, training, and operations; they enforce uncompromising safety rules regardless of warfighting urgency~\autocite{edo-4-1-5-waterfront-work-control-2025}.
  \item[\textbf{\ac{nrro}.}] Resident Naval Reactors Representative Offices monitor each nuclear-capable waterfront, enforce program standards, and can halt work if safeguards are threatened; they are independent of yard leadership to prevent production pressure from eroding safety~\autocite{edo-4-1-5-waterfront-work-control-2025}.
  \item[\textbf{\ac{rpco}.}] Reactor Plant Contractor Offices represent the prime contractors (e.g., Newport News Shipbuilding) and ensure contract compliance and technical rigor during construction, refueling, and depot maintenance~\autocite{edo-4-1-5-waterfront-work-control-2025}.
  \item[\textbf{Joint Test Groups.}] \ac{doe} and Navy personnel lead integrated test teams that certify plant reconstitution before turnover; they align with \ac{nrro}, \ac{supship}, and Ship's Force to execute rigorous testing sequences~\autocite{edo-4-1-5-waterfront-work-control-2025}.
\end{description}
Board insight: referencing the Thresher, Guitarro, and Miami investigations underscores why \ac{nrro} authority and independent QA cannot be waived, even under combat timelines.

\subsection{SUBSAFE and Submergence Systems Controls}
\begin{itemize}
  \item \textbf{\ac{subsafe}.} \ac{navsea}~0924-062-0010 mandates design, material, process, and test controls for the seawater boundary and critical systems on submarines; any work that touches a \ac{subsafe} boundary requires certified Level I material, approved procedures, and documented \ac{oqe}~\autocite{NAVSEA-0924-062-0010}.
  \item \textbf{Submergence Systems Certification.} \ac{ss800} manuals reinforce submarine rescue and diver support equipment requirements; the Submergence Systems Certification Board mirrors \ac{subsafe} rigor for non-platform systems~\autocite{SS800-AG-MAN-010}.
  \item \textbf{Level I / critical material.} \ac{navsup}/\ac{navseainst}~4440.16 controls procurement, receipt, and storage of Level I and \ac{subsafe} material to prevent counterfeit or degraded items from entering the supply chain~\autocite{NAVSEA-4440-16}.
  \item \textbf{Industrial safety.} S0570-AB-SAF-010 (fire protection) and S9002-AB-MAN-010 (submarine industrial safety) govern hot work, confined space, and damage control readiness during availabilities~\autocite{S0570-AB-SAF-010,S9002-AB-MAN-010}.
\end{itemize}

\subsection{Work Authorization and Tag-out Discipline}
\begin{enumerate}
  \item \textbf{\ac{waf}.} The \ac{nsa}/\ac{lma} generates \acp{waf} to define scope, isolations, and boundary conditions; all trades and \acp{ait} must align their work with the active \ac{waf} or request revisions through the Project Manager~\autocite{COMFLTFORCOMINST4790-3}.
  \item \textbf{Tag-out Users Manual.} S0400-AD-URM-010 prescribes red/amber tags, double barrier requirements, and control authority; the Ship's Force owns the tag-out log even when a contractor executes the work~\autocite{S0400-AD-URM-010}.
  \item \textbf{Special evolutions.} Diving, radiological work, and ordnance moves require additional approvals (e.g., Naval Sea Systems Command Supervisor of Salvage and Diving \ac{supsalv}, Naval Ordnance Safety); EDOs must ensure these are embedded in the availability plan~\autocite{NAVSEA-4740-2H,edo-4-1-5-waterfront-work-control-2025}.
\end{enumerate}

\subsection{Waterfront Organization Insights}
\begin{itemize}
  \item \textbf{Unified waterfront teams.} Naval Shipyards, Regional Maintenance Centers, \ac{rmmco}s, and Ship's Force form integrated production and test teams; weekly Project Management Reviews should cover safety, quality, schedule, and funding status~\autocite{edo-4-1-5-waterfront-work-control-2025}.
  \item \textbf{Culture of questioning attitude.} The Naval Propulsion Program teaches that any worker can stop a job; referencing the Miami (\ac{ssn}-755) fire (shipyard worker's vacuum) shows how minor lapses escalate without that culture~\autocite{edo-4-1-5-waterfront-work-control-2025}.
  \item \textbf{Documentation.} Work packages, technical work documents, and checklists must be updated in the \ac{nde} to preserve traceability for \ac{subsafe} and other certifications~\autocite{NAVSEA-0924-062-0010}.
\end{itemize}

\IfSubfilesClassLoaded{
  \RenewDocumentCommand{\entryname}{}{\textbf{\color{Modern} Acronym}}
  \RenewDocumentCommand{\descriptionname}{}{\textbf{\color{Modern} Definition}}
  \printnoidxglossary[
    type=\acronymtype,
    title=Acronyms,
    style=long-booktabs
  ]}{}
\end{document}
